Die eingeschlossene Luftmenge bleibt gleich. Bei den Bedingungen auf dem Berg würde sie nach der allgemeinen Gasgleichung $p_1V/T_1 = p_2V_2/T_2$ das Volumen
\begin{align}
 V_2 &= \VZwF \\
   &= \VF\cdot\frac{\pEi}{\pZw}\cdot\frac{\TZw}{\TEi} \\
   &= \VZw
\end{align}
einnehmen ($T_1 = \TEi$, $T_2 = \TZw$). Das ist mehr als das Volumen der Flasche: Der Druck in der Flasche ist grösser als draussen, und beim Öffnen \result{\text{strömt Luft aus}}, und zwar (bei den Bedingungen auf dem Berg gemessen)
\begin{align}
 \Delta V &= \dVFF = \dVF \approx \result{\dVFlP{-1}{2}}
\end{align}
